\begin{lemma}[Only the last PE delta remains missing among stored pairs]
Let \(u\in V(T_{q+1})\), and let
\[
 b_0<b_1<\cdots<b_p\leq q
\]
be a nonempty, strictly increasing enumeration of exactly the zero-based
rounds through \(q\) whose positive slots are PEs supported at \(u\).
Let
\[
 \mathcal S_u=\bigcup_{t=0}^{p} S_{b_t+1},
\]
where every \(S_{b_t+1}\) is the exact stored pair in that PE slot, and
let the last delta field be \(\delta_{b_p+1}=\delta_p\).  Then
\[
 \overline\varphi_{q+1}(u)\cap\mathcal S_u
 =\overline\varphi_{b_p+1}(u)\cap\mathcal S_u
 =\{\delta_p\}.
\]
The set \(\mathcal S_u\) is a finite-set indexed union, so repeated
colors are automatically identified without losing their exact values.
\end{lemma}
\begin{proof}
For each enumerated round \(b_t\), put
\[
 A_t=\overline\varphi_{b_t}(u),\qquad
 B_t=\overline\varphi_{b_t+1}(u),
\]
and denote the two colors in its PE constructor by
\(\gamma_t,\delta_t\).  We first establish
\[
 A_t\setminus\{\gamma_t\}
   =B_t\setminus\{\delta_t\}.                                  \tag{1}
\]
Indeed, eliminate the positive-slot certificate furnished by
\noderef{TashkinovSeries}.  The exact PE-kind equality leaves the PE
constructor of \noderef{TashkinovSeriesStep}.  Its
\noderef{MaximumDefectiveChoice} supplies a selected coloring
\(\pi_t\) in which the connecting edge at \(u\) has color
\(\delta_t\).  The constructor supplies a second coloring \(\rho_t\)
stable for the enlarged stored-color set.  The vertex clause of
\noderef{StableColoring} identifies the missing sets of
\(\varphi_{b_t},\pi_t,\rho_t\) at \(u\), while its edge clause fixes the
connecting \(\delta_t\)-edge under \(\rho_t\).

The rule \noderef{ParallelExtensionColor} makes \(\gamma_t\) missing at
\(u\) under these pre-switch colorings.  Thus \(\rho_t\) has a
\(\delta_t\)-edge but no \(\gamma_t\)-edge at \(u\).  By the definition
\noderef{TwoColorComponentSwap}, the PE switch on the component through
\(u\) changes the former edge to color \(\gamma_t\), and it changes no
other color at \(u\).  Consequently
\[
 B_t=(A_t\setminus\{\gamma_t\})\cup\{\delta_t\},
 \qquad
 A_t=(B_t\setminus\{\delta_t\})\cup\{\gamma_t\},                \tag{2}
\]
which proves (1).  It also shows
\(\gamma_t\notin A_t\setminus\{\gamma_t\}\) and
\(\delta_t\notin B_t\setminus\{\delta_t\}\).

No PE supported at \(u\) occurs strictly between \(b_t\) and
\(b_{t+1}\), by exactness of the enumeration.  Therefore
\noderef{HkuUnsupportedVertexColors} gives
\[
 B_t=A_{t+1}.                                                    \tag{3}
\]
The supporting-vertex field puts \(u\) in the current tree, and its
membership persists through the intervening ordered extensions by
\noderef{TreeSequenceInitialSegment} and
\noderef{TreeSequenceVertexSet}.  Meanwhile
\noderef{SupportingPEColorChain} gives
\(\delta_t=\gamma_{t+1}\).  Combining (1)--(3) inductively shows that
all \(2(p+1)\) residual sets
\[
 A_t\setminus\{\gamma_t\},\qquad
 B_t\setminus\{\delta_t\}
\]
are one common set \(R\).

The PE constructor states exactly
\(S_{b_t+1}=\{\delta_t,\gamma_t\}\).  For every \(t\), the preceding
paragraph gives \(\gamma_t\notin R\) by viewing
\(R=A_t\setminus\{\gamma_t\}\), and
\(\delta_t\notin R\) by viewing
\(R=B_t\setminus\{\delta_t\}\).  Hence the indexed union
\(\mathcal S_u\) is disjoint from \(R\).  The last instance of (2) is
\[
 B_p=R\cup\{\delta_p\}.
\]
The last PE stored set contains \(\delta_p\), so
\(B_p\cap\mathcal S_u=\{\delta_p\}\).  This proves the second displayed
intersection in the statement, including its reverse inclusion.

Finally, exactness of the enumeration says that no PE supported at
\(u\) occurs after \(b_p\) and through \(q\).  A final application of
\noderef{HkuUnsupportedVertexColors} gives
\(\overline\varphi_{q+1}(u)=B_p\); if \(b_p=q\), this is simply the same
coloring.  Intersecting with the unchanged exact finite union
\(\mathcal S_u\) proves the first displayed intersection as well.
The ambient missing sets are those of \noderef{MissingColors}, so all
set differences, unions, and intersections above are finite-set
operations exactly matching the Lean statement.
\end{proof}
